% Self-contained proofs of the retained regret lower bounds.
% Main-source antecedents: mab_aa_main.tex 1037--1179, 1215--1348.
% Original proofs: 2909--3080, 3085--3372, 3377--3569.
% The half-horizon event below replaces the original floor-threshold argument.
\section{Proofs of the retained regret lower bounds}
\label{app:retained_regret_lower_bounds}

These arguments use the original regret criterion
$R_T(\pi;\nu)=\sum_{t=1}^T(r^*(t)-r_{I_t})$ at a deterministic physical
horizon. Reward laws are stationary and independent across arms and draws.
The policies are nonanticipating, may use independent randomization, and
use the same action kernels in every environment being compared. Each
comparison keeps the deterministic arrival schedule fixed. The schedules
and reward laws may depend on $T$ where this is explicitly stated. All
logarithms are natural.

\subsection{Information inequalities with adaptive sampling}

For probability laws $P,Q$, write $D(P\Vert Q)$ for relative entropy and
$\operatorname{TV}(P,Q)=\sup_A|P(A)-Q(A)|$. For Bernoulli laws the relative
entropy is denoted by $\operatorname{kl}(p,q)$. The usual conventions apply
at zero probabilities.

\begin{lemma}[Information inequalities]
\label{lem:retained_information}
For any event $A$,
\begin{equation}
 D(P\Vert Q)\ge\operatorname{kl}(P(A),Q(A))
       \ge P(A)\log\frac1{Q(A)}-\log2.
 \label{eq:retained_binary_information}
\end{equation}
Also,
\begin{equation}
 \operatorname{TV}(P,Q)\le\sqrt{D(P\Vert Q)/2},\qquad
 |\E_P Z-\E_Q Z|\le\operatorname{TV}(P,Q)\quad(0\le Z\le1).
 \label{eq:retained_tv_information}
\end{equation}
For $0<q<1$ and $0\le p\le1$,
\begin{equation}
 \operatorname{kl}(p,q)\le\frac{(p-q)^2}{q(1-q)}.
 \label{eq:retained_bernoulli_upper}
\end{equation}
Finally, suppose two bandit environments differ only in the reward law
of arm $i$, from $P_i$ to $Q_i$, with $D(P_i\Vert Q_i)<\infty$. If their
history laws up to a deterministic time $u$ are $P^{H_u}$ and $Q^{H_u}$,
then
\begin{equation}
 D(P^{H_u}\Vert Q^{H_u})
       =\E_P N_i(u)D(P_i\Vert Q_i),\qquad
 N_i(u)=\sum_{t=1}^u\ind\{I_t=i\}.
 \label{eq:retained_adaptive_kl}
\end{equation}
\end{lemma}

\begin{proof}
Splitting the relative-entropy integral over $A$ and its complement and
applying Jensen's inequality on each part gives the first inequality in
\eqref{eq:retained_binary_information}. The remaining binary entropy
$p\log p+(1-p)\log(1-p)$ is bounded below by $-\log2$, and
$(1-p)\log(1/(1-q))\ge0$, giving the second inequality.
For fixed $q\in(0,1)$, the second derivative of
$p\mapsto\operatorname{kl}(p,q)$ is $1/(p(1-p))\ge4$; the function and its
first derivative vanish at $p=q$. Thus
$\operatorname{kl}(p,q)\ge2(p-q)^2$, with endpoint cases by limits.
Applying the first part of \eqref{eq:retained_binary_information} to
every event proves the total-variation bound in
\eqref{eq:retained_tv_information}. For bounded $Z$,
\[
 \E_P Z-\E_Q Z
       =\int_0^1\{P(Z>x)-Q(Z>x)\}\,dx,
\]
so the same definition of total variation bounds the absolute difference.
The inequality $\log x\le x-1$, applied to the two terms in Bernoulli
relative entropy, gives \eqref{eq:retained_bernoulli_upper}.

To establish \eqref{eq:retained_adaptive_kl}, include the observed rewards,
actions, and available sets in the transcript. The available sets are
deterministic and agree across the environments. Conditional on a past
history $h$, the two joint laws of the next action and reward are
$\pi_t(dj\mid h)P_j(dx)$ and $\pi_t(dj\mid h)Q_j(dx)$, respectively.
The common action kernel cancels from the likelihood ratio. Their
conditional relative entropy is therefore
$\pi_t(\{i\}\mid h)D(P_i\Vert Q_i)$, since all other laws agree.
The chain rule for relative entropy and iterated expectation give
\[
 D(P^{H_u}\Vert Q^{H_u})
   =D(P_i\Vert Q_i)\sum_{t=1}^u
           \E_P\pi_t(\{i\}\mid H_{t-1})
   =D(P_i\Vert Q_i)\E_P N_i(u).
\]
This calculation also accounts for randomized policies. The finite
horizon and finite armwise divergence ensure that its right side is
finite. These are the standard adaptive change-of-measure identities
used in bandit lower bounds \citep{kaufmann2016}.
\end{proof}

\subsection{Inferior arrivals and the static specialization}

For each $T$, let arm $0$ be available from time $1$ and uniquely optimal
in a baseline environment $\nu^{(0)}$, with mean $r^\star$. All other
means are strictly smaller. Let $\mathcal I_T$ be a nonempty finite set
of arriving arms. For $i\in\mathcal I_T$, its law is $P_i$, its mean is
$r^\star-\Delta_i$, and its arrival time is $\tau_i\le T$. An alternative
$\nu^{(i)}$ changes only this arm's law to $Q_i$, whose mean is
$r^\star+\eta_i$. Put
\[
 H_i=T-\tau_i+1,\qquad H_{\min}=\min_{i\in\mathcal I_T}H_i,\qquad
 d_i=D(P_i\Vert Q_i),\qquad
 \mathcal V_T=\{\nu^{(0)}\}\cup\{\nu^{(i)}:i\in\mathcal I_T\}.
\]
Assume $H_{\min}\to\infty$ and fixed positive constants satisfy
$\Delta_i\ge\underline\Delta$, $\eta_i\ge\underline\eta$, and
$0<d_i\le\overline d$ uniformly over the indicated arms and all
sufficiently large $T$. Let a policy sequence obey
\[
 g_T:=\sup_{\nu\in\mathcal V_T}
       \E_\nu R_T(\pi^{(T)};\nu)
       =o(H_{\min}^a)\quad\text{for every }a>0.
\]

\begin{proof}[Proof of the local-alternative regret lower bound]
Fix $i\in\mathcal I_T$ and let
\[
 A_i=\{N_i(T)\le H_i/2\}.
\]
Under the baseline, every pull of $i$ incurs regret $\Delta_i$, so
$\E_0N_i(T)\le g_T/\underline\Delta$. Markov's inequality gives
\begin{equation}
 \Prob_0(A_i^c)\le\frac{2g_T}{\underline\Delta H_i}.
 \label{eq:retained_half_baseline}
\end{equation}
Under $\nu^{(i)}$, arm $i$ is uniquely best after its arrival. Each
post-arrival pull of a different arm incurs regret at least $\eta_i$.
On $A_i$ there are at least $H_i/2$ such pulls. Therefore
\begin{equation}
 \Prob_{\nu^{(i)}}(A_i)
       \le\frac{2g_T}{\underline\eta H_i}.
 \label{eq:retained_half_alternative}
\end{equation}
Take $a=1/2$ in the uniform-goodness condition. For all sufficiently
large $T$, $g_T\le H_{\min}^{1/2}$. Consequently the right side of
\eqref{eq:retained_half_baseline} is at most
$2/(\underline\Delta H_{\min}^{1/2})$, which tends to zero uniformly in
$i$. The right side of \eqref{eq:retained_half_alternative} is at most
$2/(\underline\eta H_i^{1/2})$, which is at most $H_i^{-1/3}$ once
$H_{\min}$ is sufficiently large. Hence one common deterministic
threshold in $T$ gives, for every $i\in\mathcal I_T$,
\[
 \Prob_0(A_i)\ge\frac12,\qquad
 \Prob_{\nu^{(i)}}(A_i)\le H_i^{-1/3}.
\]
Applying \eqref{eq:retained_binary_information} and
\eqref{eq:retained_adaptive_kl} yields
\[
 d_i\E_0N_i(T)
       \ge\frac16\log H_i-\log2
       \ge\frac1{12}\log H_i
\]
for all sufficiently large $T$, uniformly over $i$. In the last step
use $H_i\ge H_{\min}\to\infty$. If an alternative event had probability
zero, the binary divergence would be infinite while the adaptive
divergence is finite; the displayed argument is valid with the usual
extended-value convention as well.

Under the same baseline environment for all the arms,
\[
 R_T(\pi^{(T)};\nu^{(0)})
       \ge\sum_{i\in\mathcal I_T}\Delta_iN_i(T).
\]
There is no summation of lower bounds obtained in incompatible baseline
environments. Taking expectations and summing the uniform armwise
bounds proves
\[
 \E_0R_T
       \ge\frac1{12}\sum_{i\in\mathcal I_T}
                  \frac{\Delta_i}{d_i}\log H_i
       \ge\frac{\underline\Delta}{12\overline d}
                 |\mathcal I_T|\log H_{\min}.
\]
\end{proof}

For a static fixed collection of $K$ inferior arms, all $H_i=T$.
Fixed positive gaps and alternative improvements give positive lower
bounds $\underline\Delta,\underline\eta$, and finitely many finite
divergences have a finite common upper bound. If additionally
$d_i\le C_i\Delta_i^2$, the preceding display becomes
\[
 \E_0R_T\ge\frac1{12}\sum_{i=1}^K
                              \frac{\log T}{C_i\Delta_i}.
\]
This proves the static specialization directly. Its fixed laws and
fixed $K$ distinguish it from a growing triangular family. The
uniform-goodness hypothesis is a condition on policies; the result
does not establish that this policy class is nonempty for every
growth schedule. In particular, its two conclusions together imply
$|\mathcal I_T|\log H_{\min}=o(H_{\min}^a)$ for every $a>0$ whenever
the assumed sequence exists.

\subsection{Successive local comparisons under benchmark improvements}

Fix $J\ge1$, positive integer block lengths $H_m$ with
$\sum_{m=1}^JH_m=T$, and endpoints $s_m=\sum_{\ell=1}^mH_\ell$,
$s_0=0$. At time $s_{m-1}+1$, two Bernoulli arms $c_m,d_m$ arrive,
with means
\[
 r_{c_m}=\theta_m,\qquad
 r_{d_m}=\theta_m+\sigma_m\varepsilon_m,
 \qquad \sigma\in\{-1,+1\}^J.
\]
The arms remain available. There are no other arms. Fix
$0<\gamma<1/2$, assume $\varepsilon_m>0$ and
$\theta_m\pm\varepsilon_m\in[\gamma,1-\gamma]$, and require
$\theta_{m+1}\ge\theta_m+\varepsilon_m$ for $m<J$.
Let $\nu_\sigma$ denote the corresponding full environment.

\begin{proof}[Proof of the deterministic-block regret lower bound]
For a fixed policy, let
\[
 \mathcal R_m=\sum_{t=s_{m-1}+1}^{s_m}
                      (r^*(t)-r_{I_t}),\qquad
 N_m=\sum_{t=s_{m-1}+1}^{s_m}\ind\{I_t=d_m\}.
\]
These are random variables, and $R_T=\sum_m\mathcal R_m$ pathwise.
Fix a block $m$ and fix all signs other than $\sigma_m$. Denote the
two resulting environments by $\nu^{m,-}$ and $\nu^{m,+}$.
For every earlier block $j<m$,
\[
 r_{c_j}=\theta_j\le\theta_m,\qquad
 r_{d_j}\le\theta_j+\varepsilon_j
                    \le\theta_{j+1}\le\theta_m.
\]
Thus the best mean in block $m$ is $\theta_m$ under the negative sign
and $\theta_m+\varepsilon_m$ under the positive sign. All arms older
than block $m$ remain available; they cannot exceed these reference
levels. Consequently,
\[
 \mathcal R_m\ge\varepsilon_mN_m
       \quad\text{under }\nu^{m,-},\qquad
 \mathcal R_m\ge\varepsilon_m(H_m-N_m)
       \quad\text{under }\nu^{m,+}.
\]
Let $P^{m,-}_{[m]}$ and $P^{m,+}_{[m]}$ be the history laws only up to
time $s_m$. The variable $N_m/H_m$ is measurable for that prefix and
lies in $[0,1]$. By \eqref{eq:retained_tv_information},
\begin{align*}
 \E_{m,-}\mathcal R_m+\E_{m,+}\mathcal R_m
 &\ge\varepsilon_m
       \{\E_{m,-}N_m+\E_{m,+}(H_m-N_m)\}\\
 &\ge H_m\varepsilon_m
       \{1-\operatorname{TV}(P^{m,-}_{[m]},P^{m,+}_{[m]})\}.
\end{align*}
Only the observations of $d_m$ have different laws. That arm has no
observations before its arrival at the beginning of block $m$.
Equation \eqref{eq:retained_adaptive_kl}, applied at the prefix
endpoint $s_m$, therefore gives
\[
 D(P^{m,-}_{[m]}\Vert P^{m,+}_{[m]})
       =\E_{m,-}N_m\,
          \operatorname{kl}(\theta_m-\varepsilon_m,
                              \theta_m+\varepsilon_m)
       \le H_m\operatorname{kl}(\theta_m-\varepsilon_m,
                                \theta_m+\varepsilon_m).
\]
Using a full-horizon transcript instead would include later pulls
of $d_m$; the prefix endpoint is essential to the bound by $H_m$.
The Bernoulli upper bound and Pinsker's inequality give
\[
 \operatorname{TV}(P^{m,-}_{[m]},P^{m,+}_{[m]})
 \le\sqrt{\frac12H_m
             \operatorname{kl}(\theta_m-\varepsilon_m,
                                 \theta_m+\varepsilon_m)}
 \le\sqrt{\frac{2H_m\varepsilon_m^2}{\gamma(1-\gamma)}}.
\]
The pair of expected block regrets is nonnegative, so it is at least
\[
 H_m\varepsilon_m
       \left(1-\sqrt{\frac{2H_m\varepsilon_m^2}
                                {\gamma(1-\gamma)}}\right)_+.
\]
Averaging over the two signs in this block contributes a factor
$1/2$. Averaging also over all fixed configurations of the other
signs, and then summing over blocks, gives
\begin{align*}
 \sup_\sigma\E_{\nu_\sigma}R_T
 &\ge2^{-J}\sum_{\sigma\in\{-1,+1\}^J}
                           \E_{\nu_\sigma}R_T\\
 &\ge\frac12\sum_{m=1}^JH_m\varepsilon_m
       \left(1-\sqrt{\frac{2H_m\varepsilon_m^2}
                                {\gamma(1-\gamma)}}\right)_+.
\end{align*}
The same uniform sign average is used for every block, so these
blockwise bounds combine under one full-environment distribution.
If $H_m\varepsilon_m^2\le\kappa\gamma(1-\gamma)$ with
$0<\kappa<1/2$, each factor is at least $1-\sqrt{2\kappa}$.
The proof holds for every policy, including randomized and
horizon-dependent policies, and therefore also after taking their
infimum.
\end{proof}

For completeness, if $\varepsilon_m\ge c_\varepsilon\sqrt{J/T}$ and
the small-information condition holds, the last bound is at least
\[
 c_\kappa\sum_{m=1}^JH_m\varepsilon_m
       \ge c_\kappa c_\varepsilon\sqrt{JT}.
\]
Conditions $H_m\asymp T/J$ and
$\varepsilon_m\asymp\sqrt{J/T}$ with a sufficiently small product
of upper constants ensure that small-information condition.
The ladder must also exist inside the bounded interval. Telescoping
the monotonicity constraints gives
\[
 \sum_{m=1}^J\varepsilon_m
       \le\theta_J+\varepsilon_J-(\theta_1-\varepsilon_1)
       \le1-2\gamma.
\]
Consequently
$c_\varepsilon J^{3/2}/\sqrt T\le1-2\gamma$, proving the necessary
range $J=O(T^{1/3})$. A homogeneous example with
$\varepsilon_m=\varepsilon$ and $\theta_m=\gamma+m\varepsilon$
is feasible whenever $(J+1)\varepsilon\le1-2\gamma$.
There are exactly $2J$ arrived arms, so block count and arm count
are different quantities. These deterministic-horizon lower bounds
do not establish matching with an upper bound at a random stopping
time, nor do they impose a growth boundary on all possible
arriving-arm environments.
